\begin{ouverture}[newton]{12}{Forces et interactions}{Mouvement et interactions}{Observer · Modéliser · Expérimenter · Expliquer}
\objectif{Dresser un bilan des actions extérieures.}
\objectif{Représenter poids, gravitation, action d’un fil ou d’un support.}
\objectif{Appliquer le principe des actions réciproques.}
\prerequis{Proportionnalité, conversions et lecture de graphiques.}
\end{ouverture}
\essentiel
\essSection{1}{Une action modélisée par une force}
Une force modélise une action exercée par un système extérieur sur le système étudié. On précise l'auteur et le receveur : $\vect F_{A\to B}$. Le vecteur porte un point d'application, une direction, un sens et une valeur en newtons.
Un diagramme d'interactions aide à recenser les acteurs : contact pour un fil ou un support ; distance pour l'attraction gravitationnelle. Il faut d'abord définir la frontière du système.
\schema[corrige]{bilan-forces}
\begin{attention}La vitesse n'est pas une force. Une « force du mouvement » ajoutée au bilan sans acteur extérieur n'est pas un modèle correct.\end{attention}
\essSection{2}{Gravitation et poids}
\begin{loi}[Gravitation]Deux masses ponctuelles, ou deux corps sphériques dont les masses sont modélisées en leur centre, s'attirent avec la valeur $F=Gm_Am_B/d^2$, où $d$ est la distance entre centres et $G=\qty{6,67e-11}{N.m^2.kg^{-2}}$. La force sur B est dirigée de B vers A.\end{loi}
Si $\vect u_{A\to B}$ va de A vers B, alors $\vect F_{A\to B}=-(Gm_Am_B/d^2)\vect u_{A\to B}$. Le signe traduit l'attraction.
\begin{definition}[Poids]Près de la surface d'une planète, $\vect P=m\vect g$. Le poids est dirigé vers le centre de la planète, approximativement vertical et vers le bas près du sol. La masse est en kg, $g$ en \unit{N.kg^{-1}}, le poids en N.\end{definition}
\essSection{3}{Actions réciproques, fil et support}
\begin{loi}[Troisième loi de Newton]$\vect F_{A\to B}=-\vect F_{B\to A}$. Les deux forces ont la même valeur et des sens opposés ; elles s'exercent sur deux systèmes différents. Elles ne s'annulent donc pas dans le bilan d'un seul système.\end{loi}
\schema[corrige]{reciprocite}
Un fil tendu exerce une traction selon sa direction, vers le fil. Un support horizontal exerce sur un objet posé au repos une action dirigée vers le haut si les autres actions horizontales sont absentes ; en présence de frottement, l'action du support peut aussi avoir une composante tangentielle.
\essSection{4}{Expérimenter et modéliser}
\begin{experience}[Mesurer une action]Accrocher un objet à un dynamomètre ; relever la force et comparer à $mg$ avec la masse mesurée. Pour observer les actions réciproques, relier deux dynamomètres et tirer doucement : comparer simultanément leurs indications. Ces indications portent sur les deux actions de l'interaction, pas sur une somme nulle appliquée à un même objet.\end{experience}
\exercices
\rubrique{Prendre confiance}
\begin{blocLivre}{ex-12-1}\begin{exercice}[1]{Nommer l’acteur}
Une main tire une valise par une poignée.
\begin{questions}
\item Quel système reçoit l’action étudiée ?
\item Comment nommer le vecteur force correspondant ?
\end{questions}
\end{exercice}\end{blocLivre}
\begin{blocLivre}{ex-12-2}\begin{exercice}[1]{Poids d’un objet}
Un objet de masse 200 g est sur Terre ; $g=\qty{10,0}{N.kg^{-1}}$.
\begin{questions}
\item Calculer son poids.
\item Préciser sa direction et son sens.
\end{questions}
\end{exercice}\end{blocLivre}
\begin{blocLivre}{ex-12-3}\begin{exercice}[1]{Contact ou distance}
Un livre repose sur une table.
\schema{bilan-forces}
\begin{questions}
\item Citer les deux actions principales sur le livre.
\item Les classer en contact ou distance.
\end{questions}
\end{exercice}\end{blocLivre}
\begin{blocLivre}{ex-12-4}\begin{exercice}[1]{Deux patineurs}
Deux patineurs se poussent mutuellement.
\schema{reciprocite}
\begin{questions}
\item Comparer les valeurs des forces qu’ils exercent l’un sur l’autre.
\item Pourquoi ces forces ne s’annulent-elles pas sur un seul patineur ?
\end{questions}
\end{exercice}\end{blocLivre}
\rubrique{S’entraîner}
\begin{blocLivre}{ex-12-5}\begin{exercice}[2]{Changer de planète}
Une masse de 3,0 kg passe de la Terre ($g=9,8$ N/kg) à la Lune ($g=1,6$ N/kg).
\begin{questions}
\item Calculer les deux poids.
\item Quelle grandeur reste la même ?
\end{questions}
\end{exercice}\end{blocLivre}
\begin{blocLivre}{ex-12-6}\begin{exercice}[2]{Une attraction gravitationnelle}
Deux masses de 10,0 kg et 20,0 kg ont leurs centres séparés de 0,50 m ; $G=\qty{6,67e-11}{N.m^2.kg^{-2}}$.
\begin{questions}
\item Calculer la valeur de l’attraction.
\item Que devient-elle si la distance double ?
\end{questions}
\end{exercice}\end{blocLivre}
\begin{blocLivre}{ex-12-7}\begin{exercice}[2]{Une boule suspendue}
Une boule de 0,200 kg est immobile au bout d’un fil vertical ; $g=10,0$ N/kg.
\schema{suspension}
\begin{questions}
\item Représenter le poids et l’action du fil.
\item Donner leurs valeurs dans cette situation d’équilibre.
\end{questions}
\end{exercice}\end{blocLivre}
\begin{blocLivre}{ex-12-8}\begin{exercice}[2]{Le centre compte}
Un satellite est à 400 km d’altitude au-dessus d’une planète de rayon 6400 km.
\begin{questions}
\item Quelle distance utiliser dans la loi de gravitation ?
\item Pourquoi 400 km ne convient-il pas ?
\end{questions}
\end{exercice}\end{blocLivre}
\rubrique{Raisonner et approfondir}
\begin{blocLivre}{ex-12-9}\begin{exercice}[3]{Terre et pomme}
Une pomme subit une attraction de 1,5 N de la Terre.
\begin{questions}
\item Quelle force la pomme exerce-t-elle sur la Terre ?
\item La grande masse de la Terre modifie-t-elle cette égalité ?
\end{questions}
\end{exercice}\end{blocLivre}
\begin{blocLivre}{ex-12-10}\begin{exercice}[3]{Un dessin à l’échelle}
Une force de 4,5 N doit être représentée à l’échelle 1 cm pour 1,5 N.
\begin{questions}
\item Calculer la longueur du vecteur.
\item Quelles informations restent nécessaires pour le tracer ?
\end{questions}
\end{exercice}\end{blocLivre}
\begin{blocLivre}{ex-12-11}\begin{exercice}[3]{Un support incliné}
Un objet est immobile sur un plan incliné grâce au frottement.
\begin{questions}
\item Pourquoi l’action totale du support n’est-elle pas nécessairement perpendiculaire au plan ?
\item Comment se dirige sa composante de frottement si l’objet tend à glisser vers le bas ?
\end{questions}
\end{exercice}\end{blocLivre}
\begin{blocLivre}{ex-12-12}\begin{exercice}[3]{Retrouver g}
À la surface d’une planète sphérique, $M=6,0\times10^{24}$ kg et $R=6,4\times10^6$ m. $G=6,67\times10^{-11}$ en SI.
\begin{questions}
\item En identifiant $mg$ à $GMm/R^2$, exprimer $g$.
\item Calculer sa valeur et commenter la dépendance à la masse de l’objet.
\end{questions}
\end{exercice}\end{blocLivre}

